\section{Limitations and discussion}\label{sec:discussion}

SA-HEACS trades a code-secrecy problem for an operations problem, and several limits remain.
\begin{itemize}
  \item \textbf{Black-box extraction.} The architecture cannot stop a patient attacker from
        cloning a function through its outputs (\cref{sec:security}). Features whose value is
        easy to learn from input--output pairs gain less from server-side execution than
        features with large hidden state or data.
  \item \textbf{Latency floor.} Interactive operations faster than about one network round
        trip cannot move to the server without prediction. Prefetch helps only when user
        behaviour is predictable, and the 40\,ms round trip assumed here is optimistic for
        mobile networks.
  \item \textbf{Offline use.} Cloud-only features fail without connectivity. The original
        blueprint's answer, locking or queueing premium features offline, is a product cost
        that users will notice.
  \item \textbf{Cost.} Remote execution moves compute from users' devices to the vendor's
        bill. Batching helps only for features with enough concurrent traffic
        (\cref{sec:eval}).
  \item \textbf{Privacy and compliance.} User content now leaves the device. Regulated data,
        data-residency rules and user trust may forbid this for some customers, and any
        server-side cache becomes a breach target.
  \item \textbf{Shared-hardware side channels.} Multi-tenant GPUs and CPUs can leak
        information across tenants. We scope this out, but a production system needs
        isolation between batches of different customers where the threat warrants it.
  \item \textbf{Simulation only.} \Cref{sec:eval} uses assumed costs and synthetic
        sessions. The numbers illustrate trade-offs, not achievable performance.
  \item \textbf{Learned control plane.} The router and prefetcher are themselves models that
        can drift or be gamed, for example by a user who crafts sequences to trigger costly
        speculation. The tier constraint is enforced in code, not learned, so a misbehaving
        model can waste money but cannot leak a cloud-only feature.
\end{itemize}

The lifecycle pipeline also changes how teams think about protection. Instead of asking
whether a feature can be protected on the client, they ask how long its secrecy is worth the
serving cost. That question has a business answer that changes over time, which is what the
promotion and demotion policy encodes.
